%vim: spl = pt
\begin{exercício}{Espectro pontual puro}{ex37}
  Um operador \(T \in \bounded(\hilbert)\) tem \emph{espectro pontual puro} ou \emph{espectro discreto puro} se \(T\) tem um conjunto ortonormal de autovetores que é total em \(\hilbert.\) Dê um exemplo de um operador auto-adjunto \(T\) com espectro pontual puro tal que \(\sigma_c(T) \neq 0.\)
\end{exercício}
\begin{proof}[Resolução]

\end{proof}
